% Part: propositional-logic
% Chapter: syntax-and-semantics
% Section: valuations-sat

\documentclass[../../../include/open-logic-section]{subfiles}

\begin{document}

\olfileid[pt]{pl}{syn}{val}

\olsection{\usetoken{P}{valuation} e Satisfação}

\begin{defn}[!!^{valuation}s]
Seja $\{\True, \False\}$ o conjunto dos dois valores de verdade,
``verdadeiro'' e ``falso''. Uma \emph{!!{valuation}} para $\Lang{L_0}$ é
uma função~$\pAssign{v}$ que atribui $\True$ ou $\False$ às
!!{propositional variable}s da linguagem, isto é, $\pAssign{v} \colon
\PVar \to \{\True, \False \}$.
\end{defn}

\begin{defn}
  Dada !!a{valuation}~$\pAssign{v}$, defina indutivamente a função de
  avaliação $\pValue{v} \colon \Frm[L_0] \to \{\True, \False \}$ por:
  \begin{align*}
    \iftag{prvFalse}{\pValue{v}(\lfalse) & = \False; \\}{}
    \iftag{prvTrue}{\pValue{v}(\ltrue) & = \True; \\}{}
    \pValue{v}(\Obj p_n) & = \pAssign{v}(\Obj p_n); \\
    \iftag{prvNot}{\pValue{v}(\lnot !A) & = \begin{cases}
        \True & \text{se } \pValue{v}(!A) = \False;\\
        \False & \text{caso contrário.}
      \end{cases} \\ }{}
    \iftag{prvAnd}{\pValue{v}(!A \land !B) & = \begin{cases}
        \True &
        \text{se $\pValue{v}(!A) = \True$ e $\pValue{v}(!B) = \True$;}\\
        \False &
        \text{se $\pValue{v}(!A) = \False$ ou $\pValue{v}(!B) = \False$}.
      \end{cases}\\}{}
    \iftag{prvOr}{\pValue{v}(!A \lor !B) & = \begin{cases}
        \True &
        \text{se $\pValue{v}(!A) = \True$ ou $\pValue{v}(!B) = \True$;}\\
        \False &
        \text{se $\pValue{v}(!A) = \False$ e $\pValue{v}(!B) = \False$}.
      \end{cases}\\}{}
    \iftag{prvIf}{\pValue{v}(!A \lif !B) & = \begin{cases}
        \True &
        \text{se $\pValue{v}(!A) = \False$ ou $\pValue{v}(!B) = \True$;}\\
        \False &
        \text{se $\pValue{v}(!A) = \True$ e $\pValue{v}(!B) = \False$}.
      \end{cases}\\}{}
    \iftag{prvIff}{\pValue{v}(!A \liff !B) & = \begin{cases}
        \True &
        \text{se $\pValue{v}(!A) = \pValue{v}(!B)$;}\\
        \False &
        \text{se $\pValue{v}(!A) \neq \pValue{v}(!B)$}.
      \end{cases}}{}
\end{align*}
\end{defn}

\begin{explain}
As cláusulas correspondem às seguintes tabelas de verdade:
\begin{center}
\iftag{prvNot}{
  \begin{tabular}{|c||c|} \hline
    $!A$ & $ \lnot !A$ \\
    \hline \hline
    $\True$ & $\False$ \\
    $\False$ & $\True$ \\
    \hline
  \end{tabular}
}{}
\iftag{prvAnd}{
  \begin{tabular}{|cc||c|} \hline
    $!A$ & $!B$ & $!A \land !B$ \\
    \hline \hline
    $\True$ & $\True$ & $\True$ \\
    $\True$ & $\False$ & $\False$ \\
    $\False$ & $\True$ & $\False$ \\
    $\False$ & $\False$ & $\False$ \\
    \hline
  \end{tabular}
}{}
\iftag{prvOr}{
  \begin{tabular}{|cc||c|} \hline
    $!A$ & $!B$ & $!A \lor !B$ \\
    \hline \hline
    $\True$ & $\True$ & $\True$ \\
    $\True$ & $\False$ & $\True$ \\
    $\False$ & $\True$ & $\True$ \\
    $\False$ & $\False$ & $\False$ \\
    \hline
  \end{tabular}
}{}%
\iftag{notprvNot,notprvAnd,notprvOr,notprvIf}{}{\\[1em]}
\iftag{prvIf}{
  \begin{tabular}{|cc||c|} \hline
    $!A$ & $!B$ & $!A \lif !B$ \\
    \hline \hline
    $\True$ & $\True$ & $\True$ \\
    $\True$ & $\False$ & $\False$ \\
    $\False$ & $\True$ & $\True$ \\
    $\False$ & $\False$ & $\True$ \\
    \hline
  \end{tabular}
}{}
\iftag{prvIff}{
  \begin{tabular}{|cc||c|} \hline
    $!A$ & $!B$ & $!A \liff !B$ \\
    \hline \hline
    $\True$ & $\True$ & $\True$ \\
    $\True$ & $\False$ & $\False$ \\
    $\False$ & $\True$ & $\False$ \\
    $\False$ & $\False$ & $\True$ \\
    \hline
  \end{tabular}
}{}
\end{center}
\end{explain}

\begin{prob}
Considere acrescentar a $\Lang{L_0}$ um conectivo ternário $\diamondsuit$
com avaliação dada por
\begin{gather*}
  \pValue{v}(\diamondsuit( !A , !B , !C ) ) = \begin{cases}
    \pValue{v}( !B ) &
    \text{se $\pValue{v}(!A) = \True$;}\\
    \pValue{v}( !C ) &
    \text{se $\pValue{v}(!A) = \False $}.
  \end{cases}
\end{gather*}
Escreva a tabela de verdade desse conectivo.
\end{prob}

\begin{thm}[Determinação Local]
\ollabel{thm:LocalDetermination} Suponha que $\pAssign{v_1}$ e
$\pAssign{v_2}$ sejam !!{valuation}s que coincidam nas !!{propositional
variable}s que ocorrem em $!A$, isto é, $\pAssign{v_1}(\Obj p_n) =
\pAssign{v_2}(\Obj p_n)$ sempre que $\Obj p_n$ ocorrer em alguma
!!{formula}~$!A$. Então $\pValue{v_1}$ e $\pValue{v_2}$ também coincidem
em~$!A$, isto é, $\pValue{v_1}(!A) = \pValue{v_2}(!A)$.
\end{thm}

\begin{proof}
Por indução sobre $!A$.
\end{proof}

\begin{defn}[Satisfação]
\ollabel{defn:satisfaction} Podemos definir indutivamente a noção de
  \emph{satisfação de !!a{formula}~$!A$ por
  !!a{valuation}~$\pAssign{v}$}, $\pSat{v}{!A}$, da seguinte maneira.
  (Escrevemos $\pSat/{v}{!A}$ com o significado de ``não $\pSat{v}{!A}$.'')
\begin{enumerate}
\tagitem{prvFalse}{%
  \indcase{!A}{\lfalse}{$\pSat/{v}{\indfrm}$.}}{}

\tagitem{prvTrue}{%
  \indcase{!A}{\ltrue}{$\pSat{v}{\indfrm}$.}}{}

\item \indcase{!A}{\Obj p_i}{$\pSat{v}{\indfrm}$
  sse $\pAssign{v}(\Obj p_i) = \True$.}

\tagitem{prvNot}{%
  \indcase{!A}{\lnot !B}{$\pSat{v}{\indfrm}$ sse
    $\pSat/{v}{!B}$.}}{}

\tagitem{prvAnd}{%
  \indcase{!A}{(!B \land !C)}{$\pSat{v}{\indfrm}$ sse $\pSat{v}{!B}$
    e $\pSat{v}{!C}$.}}{}

\tagitem{prvOr}{%
  \indcase{!A}{(!B \lor !C)}{$\pSat{v}{\indfrm}$ sse
    $\pSat{v}{!B}$ ou $\pSat{v}{!C}$ (ou ambos).}}{}

\tagitem{prvIf}{%
  \indcase{!A}{(!B \lif !C)}{$\pSat{v}{\indfrm}$ sse $\pSat/{v}{!B}$
    ou $\pSat{v}{!C}$ (ou ambos).}}{}

\tagitem{prvIff}{%
  \indcase{!A}{(!B \liff !C)}{$\pSat{v}{\indfrm}$ sse ambos
    $\pSat{v}{!B}$ e $\pSat{v}{!C}$, ou nem $\pSat{v}{!B}$
    nem $\pSat{v}{!C}$.}}{}
\end{enumerate}
Se $\Gamma$ é um conjunto de !!{formula}s, $\pSat{v}{\Gamma}$ sse
$\pSat{v}{!A}$ para todo~$!A \in \Gamma$.
\end{defn}

\begin{prop}\ollabel{prop:sat-value}
  $\pSat{v}{!A}$ sse $\pValue{v}(!A) = \True$.
\end{prop}

\begin{proof}
  Por indução sobre~$!A$.
\end{proof}

\begin{prob}
  Demonstre \olref[pl][syn][val]{prop:sat-value}
\end{prob}

\end{document}
